% Vietnamese translation for OI Public Library
% Source-PDF: IOI中国国家候选队论文/2001/李益明.pdf
\documentclass[11pt]{article}
\usepackage{oipl-vn}

\title{Một Số Khảo Sát Về Bài Toán Hình Học}
\author{Li Yiming}
\date{2001}

\begin{document}
\maketitle

\origtitle{Speech draft on computational geometry}
\sourcepath{IOI Candidate Team Papers/2001/Li Yiming.pdf}

\begin{abstract}
Bài viết là một bản phát biểu ngắn về cách tiếp cận bài toán hình học trong
thi tin học. Tác giả nhấn mạnh vai trò của các thuật toán hình học cơ bản,
những thuật toán kinh điển, và đặc biệt là việc chuyển hóa mô hình hình học
thành mô hình hiệu quả hơn như đồ thị hoặc xấp xỉ.
\end{abstract}

\section{Mở đầu}

Hình học tính toán nghiên cứu các thuật toán cho bài toán hình học. Trong kỹ
thuật và toán học hiện đại, các lĩnh vực như đồ họa máy tính, thiết kế hỗ trợ
máy tính và robot học đều dùng hình học tính toán. Trong thi tin học, bài hình
học cũng đã bắt đầu xuất hiện, nhưng trong thực tế thi đấu, tỷ lệ điểm của các
bài hình học thường là thấp nhất. Vì vậy tác giả thử khảo sát một số thuật
toán và kỹ thuật cho loại bài này.

Bất kỳ thuật toán phức tạp nào cũng được ghép từ nhiều thuật toán đơn giản.
Bài hình học cũng vậy. Trước hết cần nắm vững những thao tác cơ bản:
\begin{enumerate}
  \item tính hệ số góc của đường thẳng;
  \item tìm giao điểm của hai đường thẳng;
  \item kiểm tra hai đoạn thẳng có cắt nhau hay không;
  \item tính tích có hướng, v.v.
\end{enumerate}
Đó là nền tảng của bài hình học. Không quen thuộc với các thao tác cơ bản này
rất dễ dẫn đến thất bại.

Tuy nhiên chỉ có thao tác cơ bản vẫn chưa đủ. Nếu trong phòng thi mới tạm thời
nghĩ cách ghép chúng lại thì thường không kịp. Ta còn cần quen với một số
thuật toán kinh điển để có thể dùng trực tiếp, chẳng hạn:
\begin{enumerate}
  \item tìm bao lồi;
  \item tìm cặp điểm gần nhất;
  \item kiểm tra một điểm có nằm trong đa giác hay không.
\end{enumerate}

\section{Một số loại bài hình học}

\subsection{Bài tính toán thuần túy}

Loại bài này cần nền tảng hình học giải tích vững chắc, đồng thời phải nhìn
vấn đề toàn diện và phân tích kỹ các trường hợp đặc biệt. Ví dụ, khi tính hệ
số góc cần xét đường thẳng đứng; khi tìm giao điểm hai đường thẳng cần xét
trường hợp song song. Những chi tiết như vậy chủ yếu dựa vào tích lũy trong
quá trình học tập.

\subsection{Bài tồn tại}

Loại bài này có thể giải trực tiếp bằng tính toán: tìm được nghiệm khả thi thì
tồn tại, không tìm được thì không tồn tại. Nhưng hiệu quả của mô hình phụ
thuộc vào mức độ trừu tượng hóa. Mô hình càng trừu tượng thì càng hiệu quả.
Mô hình hình học thường có mức trừu tượng thấp; hơn nữa bài tồn tại thường có
nhiều bộ dữ liệu trong một test, nên mô hình hình học trực tiếp thường không
đủ nhanh. Vì vậy cần biến đổi mô hình hình học thành một mô hình hiệu quả hơn.

\subsection{Bài tối ưu hình học}

Đây là loại khó trong các bài hình học. Nói chung không có thuật toán phổ quát
nào thật hiệu quả để tìm nghiệm tối ưu, nên cách thường dùng là thuật toán xấp
xỉ. Chất lượng thuật toán xấp xỉ phụ thuộc hoàn toàn vào khoảng cách giữa
nghiệm tìm được và nghiệm tối ưu.

\section{Ví dụ 1: Hotter--Colder}

Người chơi \(B\) chọn một điểm mục tiêu trên tờ giấy \(100\times100\). Ban đầu
người chơi \(A\) đứng tại \((0,0)\). Mỗi lần \(A\) chuyển từ điểm hiện tại đến
một điểm mới. Nếu điểm mới gần mục tiêu hơn, \(B\) nói \texttt{Hotter}; nếu xa
hơn, \(B\) nói \texttt{Colder}; nếu khoảng cách không đổi, \(B\) nói
\texttt{Same}. Dữ liệu vào gồm nhiều dòng, mỗi dòng chứa điểm \((x,y)\) mà
\(A\) vừa tới và câu trả lời của \(B\). Với mỗi câu trả lời, cần tính diện
tích miền có thể chứa điểm mục tiêu, chính xác đến hai chữ số sau dấu phẩy.

Đây là một bài tính toán thuần túy. Trước hết chứng minh rằng miền khả dĩ luôn
là một đa giác lồi. Mỗi câu trả lời của \(B\) xác định mục tiêu nằm ở một phía
của đường trung trực giữa điểm xuất phát và điểm đến, hoặc nằm ngay trên đường
trung trực ấy. Vì mỗi lần cắt bởi một nửa mặt phẳng, miền khả dĩ vẫn là đa
giác lồi; giao của nhiều đa giác lồi cũng là đa giác lồi.

Cách giải là khởi tạo miền khả dĩ bằng hình vuông
\[
  (0,0),(0,100),(100,100),(100,0).
\]
Sau mỗi câu trả lời của \(B\), dùng đường trung trực tương ứng để cắt đa giác,
rồi giữ phần còn có thể chứa mục tiêu.

Tuy nhiên cách làm này chỉ đúng khi xử lý đủ các trường hợp đặc biệt. Ví dụ,
nếu điểm mới trùng với điểm cũ thì không tồn tại đường trung trực; đây là một
điểm cần đặc biệt chú ý. Bài toán dùng nhiều kiến thức hình học giải tích:
tìm giao điểm đoạn thẳng, xác định hai điểm nằm ở hai phía của một đường
thẳng, chứng minh miền cuối cùng là đa giác lồi, và tính diện tích đa giác.

\section{Ví dụ 2: Đường đi qua các cọc}

Trên một dải đường dài vô hạn có \(n\) cọc, với \(n\le200\), thể tích cọc xem
như bằng \(0\). Một người muốn đi từ bên trái sang bên phải. Người đó được
xem gần đúng là một hình tròn bán kính \(R\). Hỏi có thể đi qua hay không?

Cách cơ bản nhất là quét từng cột thẳng đứng từ cọc trái nhất tới cọc phải
nhất, tính các khoảng có thể đi được. Nếu tại cột chứa cọc phải nhất vẫn còn
khoảng đi được thì có lời giải, ngược lại thì không. Nhưng nếu hai cọc ngoài
cùng cách nhau rất xa thì độ phức tạp của phép quét này rất cao, nên cần
chuyển hóa mô hình.

Trong hình ban đầu, các cọc đứng yên và người di chuyển như một hình tròn; xử
lý trực tiếp như vậy khá phiền. Ta chuyển người thành một điểm, cụ thể là tâm
hình tròn. Khi đó tâm người phải cách mọi cọc ít nhất \(R\), nên mỗi cọc được
thay bằng một hình tròn tâm tại cọc, bán kính \(R\), còn người chỉ còn là một
điểm. Việc tính miền đi được trở nên dễ hơn.

Nhưng mô hình này vẫn chưa giải quyết triệt để bài toán. Hai thuật toán trước
đều tính phần bên ngoài các hình tròn, trong khi xét tính liên thông của phần
bên trong các hình tròn lại đơn giản hơn. Vì dải đường vô hạn theo phương
trái--phải, nếu hai phần trái và phải liên thông trong miền ngoài các hình
tròn thì hai biên trên và dưới không liên thông trong miền bên trong các hình
tròn. Ngược lại, nếu trái và phải không liên thông ở miền ngoài, thì trên và
dưới liên thông trong miền bên trong.

Do đó ta chuyển từ quét theo từng cột sang xét liên thông của các hình tròn.
Quan sát tính chất của hình tròn: hai điểm bất kỳ trong cùng một hình tròn
đều liên thông. Vì các hình tròn có cùng bán kính, không có quan hệ chứa nhau
phức tạp; giữa hai hình tròn chỉ cần xét giao nhau hay tách rời. Nếu hai hình
tròn giao nhau, hai miền cấm tương ứng liên thông, nên nối một cạnh giữa hai
tâm. Thêm một đỉnh nguồn nối với mọi hình tròn cắt biên trên, và một đỉnh
đích nối với mọi hình tròn cắt biên dưới. Bài toán hình học được chuyển hoàn
toàn thành bài toán đồ thị: chỉ cần kiểm tra nguồn và đích có đường đi hay
không.

Từ quá trình trên có thể rút ra kết luận: với loại bài hình học này, cần hiểu
rõ các phép biến đổi hình học, khai thác các manh mối ẩn trong đề, và tìm ra
bản chất của bài toán.

\section{Ví dụ 3: Đặt bò trong hình chữ nhật}

Một nông dân thả \(n\) con bò trên một thửa ruộng hình chữ nhật kích thước
\(x\times y\), với \(n\le25\). Các con bò rất ghét nhau, nên chúng muốn ở càng
xa nhau càng tốt. Tiêu chuẩn của chúng là tổng nghịch đảo khoảng cách tới các
con bò khác càng nhỏ càng tốt. Người nông dân muốn làm chúng vui nhất có thể,
nên cần tìm cách đặt các con bò sao cho tổng ấy nhỏ nhất. Bài chấm theo mức
độ gần với nghiệm tối ưu. Dữ liệu vào là \(x,y,n\), dữ liệu ra là vị trí của
\(n\) con bò.

Rõ ràng đây là bài tối ưu hình học và nên dùng thuật toán xấp xỉ. Chất lượng
lời giải phụ thuộc vào chất lượng thuật toán xấp xỉ.

Cách đơn giản nhất là dùng nghiệm đẹp đã tính tay trên hình vuông rồi co giãn
theo tỷ lệ sang hình chữ nhật. Chẳng hạn với \(n=4\), nghiệm trên hình vuông
là bốn góc; khi \(x\) và \(y\) không chênh lệch nhiều, đây cũng là nghiệm tối
ưu hoặc gần tối ưu trên hình chữ nhật. Nhưng khi \(x\) và \(y\) chênh nhau
lớn, cách này kém xa nghiệm tối ưu, gần như không được điểm.

Trong bài tối ưu chính xác, tham lam thường khó phát huy tác dụng; nhưng nó
lại là một thuật toán xấp xỉ rất thường dùng. Với bài này có thể thử tham
lam. Chọn điểm đầu là \((0,0)\). Mỗi điểm tiếp theo được chọn sao cho tổng
nghịch đảo khoảng cách tới các điểm đã đặt là nhỏ nhất. Khi tìm điểm tốt nhất
cần dùng phương pháp tinh chỉnh dần. Cách này cho nghiệm khá gần tối ưu, và
khó tìm dữ liệu khiến tham lam quá xa nghiệm tối ưu; vì vậy trong bài này
tham lam là một thuật toán hiệu quả.

Dĩ nhiên tham lam vẫn có khoảng cách với nghiệm tối ưu và khó đạt toàn bộ
điểm. Vì bài không có lời giải cố định tối ưu duy nhất, mọi thuật toán cố định
đều có thể gặp dữ liệu khiến nó không đạt điểm tối đa. Do đó có thể nghĩ đến
thuật toán không cố định, tức thuật toán ngẫu nhiên. Bản gốc chỉ nêu gợi ý này
và để người đọc tự suy nghĩ thêm.

\section{Phụ lục: diễn giải các chương trình}

PDF gốc kèm ba chương trình Pascal. Dưới đây là bản diễn giải thuật toán bằng
tiếng Việt, giữ lại ý chính của chương trình mà không lặp lại các chú thích
tiếng Trung trong mã nguồn.

\subsection{Bài thứ nhất: đường qua các cọc}

Đọc chiều rộng dải đường, bán kính \(R\), số cọc và tọa độ các cọc. Dựng một
đồ thị vô hướng có các đỉnh \(1..n\) ứng với các cọc, thêm đỉnh \(0\) biểu
diễn biên trên và đỉnh \(n+1\) biểu diễn biên dưới. Nếu một cọc cách biên trên
không quá \(R\), nối nó với đỉnh \(0\); nếu cách biên dưới không quá \(R\),
nối nó với đỉnh \(n+1\). Với hai cọc \(i,j\), nếu
\[
  (x_i-x_j)^2+(y_i-y_j)^2 \le (2R)^2,
\]
thì hai hình tròn bán kính \(R\) giao nhau, nên nối cạnh \(i\)--\(j\).
Sau đó dùng Floyd--Warshall hoặc tìm kiếm đồ thị để kiểm tra \(0\) có nối tới
\(n+1\) hay không. Nếu có, các miền cấm nối liền hai biên và người không thể
đi qua; nếu không, có thể đi qua.

\subsection{Bài thứ hai: Hotter--Colder}

Khởi tạo đa giác khả dĩ là hình vuông \(100\times100\), và vị trí trước của
người chơi \(A\) là \((0,0)\). Với mỗi bước mới \((x,y)\), nếu câu trả lời là
\texttt{Same} thì miền khả dĩ bị thu xuống đường trung trực; chương trình gốc
xử lý bằng cách đặt số đỉnh về \(0\). Nếu điểm mới trùng điểm cũ thì bỏ qua
bước này.

Ngược lại, lập phương trình đường trung trực giữa \((x_p,y_p)\) và \((x,y)\):
\[
  2(x-x_p)X+2(y-y_p)Y+x_p^2+y_p^2-x^2-y^2=0.
\]
Dựa vào \texttt{Hotter} hoặc \texttt{Colder}, chọn nửa mặt phẳng đúng. Với
mỗi cạnh của đa giác hiện tại, nếu hai đầu nằm hai phía khác nhau của đường
thẳng thì chèn giao điểm với đường trung trực. Sau đó xóa các đỉnh nằm ở nửa
mặt phẳng sai và tính diện tích đa giác còn lại bằng cách chia nó thành các
tam giác cùng chung một đỉnh.

\subsection{Bài thứ ba: đặt bò bằng tham lam}

Đặt con bò thứ nhất ở \((0,0)\); nếu có con thứ hai thì đặt ở \((x,y)\). Với
mỗi con bò tiếp theo, bắt đầu từ tâm hình chữ nhật và dùng tìm kiếm tinh chỉnh
dần trên một lưới \(5\times5\) quanh điểm hiện tại. Với mỗi điểm ứng viên
trong hình chữ nhật và không trùng điểm đã dùng, tính
\[
  \sum_{\text{bò đã đặt }p}\frac{1}{\operatorname{dist}(p,q)}.
\]
Chọn điểm có tổng nhỏ nhất, dịch tâm tìm kiếm về điểm đó, rồi giảm bước lưới
theo tỷ lệ. Khi bước đủ nhỏ, ghi nhận điểm tìm được làm vị trí của con bò tiếp
theo. Cuối cùng in toàn bộ tọa độ với hai chữ số thập phân.

\end{document}
